% Section 9 -- Related work
\section{Related work}
\label{sec:related}

Three lines converge on joint BGP--traffic attribution: routing anomaly
detection, the statistical and counterfactual machinery for causal
attribution, and the selective-prediction discipline that governs when a
system should refuse to answer.

\subsection{Routing anomaly detection}

Early detection worked entirely in the control plane, comparing
announcements against an expectation of who should be announcing a
prefix~\cite{rfc6811}. ARTEMIS operationalised this at network scale by
comparing monitored BGP updates with prefix-origin expectations and
alerting within a minute~\cite{artemis}. The limitation is structural:
control-plane observation cannot separate an announcement that degraded
reachability from one that was merely unusual, and it emits an unordered
candidate stream that still requires operator adjudication. RPKI origin
validation addressed the cryptographic part of the same
question~\cite{rfc6811,rodday-2024-rpki}, and measurements quantify how
much risk survives partial deployment~\cite{rpki-coming-of-age-2019} and
how origin validation is used in operational practice~\cite{rovista-2023}.

A second line moved to the data plane. Probes and distributed vantage
points infer reverse paths and answer whether connectivity actually
broke~\cite{netalyzr-2013}; their coverage is bounded by vantage-point
count and their cost grows with address space, so they complement rather
than replace control-plane monitoring. Learned detectors then fused the
two, scoring prefixes from AS-path and update-stream features or from the
semantics of routing messages~\cite{beam2024}. These improve per-event
precision but inherit the supervision problem: labelled incidents are
scarce, and their outputs are scores rather than auditable arguments.
Leak-specific systems turn the valley-free violation into an enforcement
mechanism between consenting ASes~\cite{peerlock}, although the problem
definition itself is standardised~\cite{rfc7908,rfc9234}; because
participation is voluntary and the mechanism protects one leak class, they
reduce incidence without providing attribution. Our results add a reason
why that division is durable: on our data the leak class is exactly where
the data-plane channels carry least information
(\S\ref{sec:eval-direction}).

\subsection{Causal attribution from telemetry}

Granger asked whether the past of one series improves prediction of the
present of another~\cite{granger-1969-investigating}. It is the natural
starting point for BGP--traffic analysis and, in our measurements, the
weakest: the test wants stationary residuals, and update streams are
bursty by construction, so its precondition fails. Transfer entropy
relaxed linearity by working in the information domain~\cite{schreiber-2000-te,barnett-2009-granger-network}
and, in our evaluation, is the best-performing single channel on
$T\to R$ events. It inherits a second problem: the estimate depends on
embedding dimension and history length, and it yields a scalar with no
notion of when the evidence is too weak to support a claim. We keep both
estimators and add what they lack---a surrogate null and a floor.

The counterfactual literature asks a strictly stronger question.
Pearl's probability of necessity and sufficiency asks whether treatment
was necessary for the effect and sufficient to produce
it~\cite{pearl-causality-2009,pearl1999}; identification results were
sharpened by Shpitser and Pearl and extended to graph-free
settings~\cite{shpitser-2008-pns,pearl-2014-pns-3}. The potential-outcomes
framing makes the same point from the other direction: a causal contrast
is only defined against a treatment assignment the analyst can
defend~\cite{rubin-2005-causal}, and mediation analysis shows how much of
an observed effect a specified mechanism explains once that assignment
exists~\cite{imai-2010-mediation}. The gap for network telemetry is that
no a priori causal graph exists: BGP updates and traffic are jointly
produced by the same operator decisions, so the graph the theorem requires
is exactly what the analyst lacks, and graph-free estimators reduce to
difference-in-means on a pre/post partition---the estimator we use.

Recent network-operations work imports this machinery directly, learning
causal graphs or counterfactual contrasts over routing and traffic
telemetry. We are not aware of a published system in this line that
reports an abstention record alongside its directional accuracy, so our
comparison is with the individual channels rather than with a published
estimator (\S\ref{sec:eval-direction}), and we require each channel's
significance condition to hold before the engine commits. Given that our
own point estimates produced two wrong commitments out of two on the
real-cache set, we read the point-estimate style of reporting as the main
open risk in the area rather than as a baseline we beat.

\subsection{Selective prediction and auditability}

Refusing to answer is a classical option: Chow derived the optimal
error--reject tradeoff for a fixed classifier~\cite{chow-1970-rejection},
and selective classification makes the abstention threshold a
first-class design choice~\cite{geifman-2017-selective}. We borrow the
framing but not the guarantee. Our veto conjunction is a decision
procedure with stated per-veto tests, not an optimal-risk construction,
and the coverage we report is what that literature would call the
acceptance rate---low, and reported as such.

The dependability community frames attribution as a fault-tolerance
problem with reproducibility as a review criterion. Because evaluation
sets here are necessarily small, interval estimates and permutation
tests remain the workhorse tools~\cite{phipson-2010-permutation,efron-1979-bootstrap},
and a proportion measured on few trials is conventionally summarised by a
Wilson interval~\cite{wilson-1927}. We follow that practice: every number
in \S\ref{sec:eval} traces to a committed artefact.

\subsection{Positioning}

No prior system provides a rule for combining control-plane and
data-plane evidence into a single adjudicated verdict with an explicit
refusal condition. Phase~1 produces a five-class correlation verdict
that tells an operator whether a routing anomaly has traffic impact;
Phase~2 attempts directional attribution and declines when the channels
do not clear their nulls. We do not claim to be first at joint
attribution, and we do not claim state-of-the-art direction accuracy,
which on our data no method achieves. The claim is narrower: a
deterministic, replayable procedure whose failure mode is a documented
abstention rather than a confident error.

On the agent side we do not claim that multi-agent decomposition improves
verdicts on every distribution (\S\ref{sec:agent-monolithic}). The
literature we draw on is concerned with the two properties we do rely
on---that an agent's intermediate state be inspectable, and that a failed
component not corrupt the decision---and what we add is a measurement of
both.
